\chapter{Reference manual for records and fixtures}
\label{ch:reference-manual}

The preceding appendix explains why the records exist. This one makes them
usable. An engineer should be able to build a first fixture from these pages;
a researcher should be able to trace a load-bearing sentence to its source; and
a sponsor should be able to inspect a decision packet without opening a
repository log. The manual is intentionally concrete. It belongs after the
argument because its tables and field names are tools for reproduction, not the
story the reader must learn first.

\section{Document and version fields}
\label{sec:manual-document-fields}

A document record begins with the work, not with a model run. The minimum fields
are shown below.

\begin{table}[H]
\centering
\small
\begin{tabularx}{\textwidth}{@{}p{0.23\textwidth}p{0.24\textwidth}L@{}}
\toprule
Field & Type & Interpretation \\
\midrule
document\_id & stable string & Identifies the work across revisions and
renderings \\
genre & controlled text plus free description & Research proposal, model note,
memorandum, survey, or another declared class \\
reader\_role & free text & The expertise and decision responsibility of the
intended reader \\
decision\_question & prose & The action or judgment the document is meant to
support \\
exemplars & ordered list of source ids & Documents that define the intended
explanatory pace for this reading task \\
known\_vocabulary & list of terms & Terms the reader may reasonably know
without a first-use gloss \\
exemption\_zones & line ranges plus reason & Preview, map, decision-box, or
reference regions where names may precede teaching \\
brief\_version & integer plus date & Cumulative reader directives used by this
inspection \\
source\_root & path under the authorized boundary & Directory whose contents
are included in the snapshot \\
parent\_version & nullable identifier & Previous source snapshot, if the
document is being revised \\
assistance\_disclosure & structured record & Tools, model use, human editing,
and unresolved disclosure choices \\
privacy\_class & controlled text & Public, internal, confidential, or restricted \\
\bottomrule
\end{tabularx}
\caption{Fields that define the document before diagnostics run.}
\label{tab:manual-document-fields}
\end{table}

The free description beside a controlled genre is important. Two documents may
both be called a proposal while one asks for a research grant and the other asks
for authorization to change a production model. The reader task, not the label,
determines which evidence and wording are appropriate. A missing decision
question is an intake failure, not an invitation to let a model infer one.

A version record adds build information and source lineage:

\begin{table}[H]
\centering
\small
\begin{tabularx}{\textwidth}{@{}p{0.23\textwidth}p{0.24\textwidth}L@{}}
\toprule
Field & Type & Interpretation \\
\midrule
version\_id & immutable string & One source and configuration snapshot \\
entry\_point & path & Top-level TeX, Markdown, or other source file \\
route\_files & ordered list & Included files and their content hashes \\
build\_id & string & Compiler and dependency environment \\
source\_sha256 & hexadecimal string & Exact input identity \\
rendered\_sha256 & hexadecimal string & Exact PDF identity \\
protected\_manifest & path or embedded object & Objects extracted for comparison \\
build\_result & pass, fail, or incomplete & Typography and reproducibility result,
not a reader judgment \\
\bottomrule
\end{tabularx}
\caption{A version record ties source, build, and rendered output together.}
\label{tab:manual-version-fields}
\end{table}

The distinction between a source hash and a rendered hash matters when TeX
changes a page break without changing a sentence, or when a source edit is
present but the PDF was not rebuilt. Both conditions should be visible in a
comparison packet.

\section{Protected-object dictionary}
\label{sec:manual-object-fields}

Each protected object has a typed identity. The following fields are sufficient
for the first release:

\begin{table}[H]
\centering
\small
\begin{tabularx}{\textwidth}{@{}p{0.21\textwidth}p{0.27\textwidth}L@{}}
\toprule
Object type & Required fields & Normalization rule \\
\midrule
Equation & source span, environment, label, variables, units, rendered anchor &
Normalize whitespace and macro aliases; retain signs, exponents, limits, and
unknown commands \\
Number & value, unit, precision, context, source span & Normalize separators and
formatting; do not round or drop a sign \\
Citation & key, authors, title, identifier, sentence span & Normalize key and
identifier; preserve sentence and bibliography versions \\
Table & caption, row/column labels, cells, units, source span & Normalize cell
formatting; retain denominator, ordering, and missing status \\
Qualification & text span, associated finding or result, scope terms & Preserve
scope words and link to the result they qualify \\
Label or cross-reference & key, target, source span, rendered page & Compare
key, target, and resolved page independently \\
Quotation or code & delimiters, source, language, text hash & Preserve content
exactly unless an author-approved replacement is recorded \\
\bottomrule
\end{tabularx}
\caption{Typed protected objects keep normalization from becoming silent
rewriting.}
\label{tab:manual-object-dictionary}
\end{table}

A normalization function is a convenience for matching, not a replacement for
the source. Let $g_k$ be the function for object kind $k$. The comparison
stores both $raw(p)$ and $g_k(raw(p))$:
\begin{equation}
  record(p)=\bigl(kind(p),raw(p),g_{kind(p)}(raw(p)),loc(p),owner(p)\bigr).
  \label{eq:manual-object-record}
\end{equation}
When a later implementation changes $g_k$, it must replay old fixtures and
retain the prior normalized value. Otherwise a dependency update can make a
historical change appear to have been unchanged.

A qualification deserves special treatment. It is not merely a sentence with
words such as "may" or "under." The record links it to the proposition whose
scope it limits. If the proposition moves, the system asks the author to
confirm the link. This catches the common repair in which a caveat survives as
a grammatical sentence but no longer qualifies the result it was written to
bound.
\section{The fixture catalogue}
\label{sec:manual-fixtures}

The first benchmark corpus should contain small fixtures that isolate one
failure and a few composite documents that reproduce the reader problem. Each
fixture has a positive behavior, a negative behavior, and an abstention rule.

\begin{description}[style=nextline,leftmargin=33mm,labelindent=0mm]
\item[Macro fixture] A custom command expands to a number in one table and a
label in another. A parser must preserve the call and expansion or abstain.
\item[Equation fixture] The child changes an exponent and then changes only
whitespace. The first difference must be substantive; the second may be
formatting under the declared normalization.
\item[Table fixture] A percentage is recomputed with a different denominator.
The system must identify the denominator and show the before-and-after values.
\item[Citation fixture] A key resolves to a real paper whose method does not
support the sentence. Identity resolution passes; support review remains open.
\item[Qualification fixture] A sentence moves a population or horizon
qualification away from the result it limits. The object comparison must keep
the association visible.
\item[Register fixture] A private phase label and repository path occur in the
opening paragraph. The finding should point to the public decision the reader
needs, not merely mark a forbidden word.
\item[Intentional-pattern fixture] Three adjacent sentences have the same form
because they enumerate the steps of a derivation. A style rule should either
abstain or give the author a reason to retain the pattern.
\item[Pacing fixture] Three prose paragraphs introduce a known set of terms at
different rates. The meter must report the burst locations and compare the
target with supplied exemplars without declaring a universal pass value.
\item[First-use fixture] An acronym appears before its expansion, followed by
the same acronym after a generic gloss. The audit flags the first occurrence,
accepts the second, and respects a declared preview exemption.
\item[Assertion fixture] A prose repair adds a number or superlative without a
citation. The comparison opens a question and never silently accepts the new
claim.
\item[Reader fixture] Two versions have identical equations but different
openings. A reader unfamiliar with the project should be able to identify which version makes
the decision and mechanism clearer.
\end{description}

The composite corpus then combines these cases in realistic documents. It
should include at least one proposal, one mathematical note, and one
investment-style memorandum, with a source snapshot and a human-authored answer
key for each. The answer key records what a competent reader should be able to
recover, not which recommendation the reader must endorse.

For fixture $i$, the expected result is a partial function
\begin{equation}
  E_i: input_i \longrightarrow
  \{\mathrm{pass},\mathrm{fail},\mathrm{abstain}\}.
  \label{eq:manual-fixture-result}
\end{equation}
A candidate that returns "pass" on an unresolved negative fixture has a false
pass even if its aggregate accuracy is high. The benchmark report lists every
fixture, including skipped or unsupported cases, so the denominator cannot
quietly improve when difficult documents are removed.

\section{Finding record and author decision}
\label{sec:manual-finding-record}

A finding is useful when another person can understand it without opening the
tool's implementation. The recommended fields are:

\begin{table}[H]
\centering
\small
\begin{tabularx}{\textwidth}{@{}p{0.24\textwidth}p{0.25\textwidth}L@{}}
\toprule
Field & Example & Purpose \\
\midrule
finding\_id & F-014 & Stable reference in a revision packet \\
location & chapter.tex:87, PDF page 6 & Lets the author inspect source and
rendered context \\
observation & Private label precedes public definition & States what is visible \\
consequence & New reader may not recover the comparison & Explains why attention
might be worthwhile \\
question & Name the baseline and changed mechanism & Gives the author a usable
next move \\
suggestion & Optional text or none & Offers an alternative without authority \\
priority & provisional ordinal & Orders attention; never certifies quality \\
author\_decision & accept, reject, defer, rewrite & Preserves human ownership \\
decision\_reason & short prose & Explains a rejection or substantive change \\
unit & chapter or section id & Keeps pacing and ordering work bounded and
replayable \\
repair\_move & concrete-first, unfold, read-back, generic-first-name-after,
or other named move & Gives the author a finite option without prescribing
wording \\
detector\_evidence & metric, exemplar ids, or first-use signal & Shows why the
question was raised and what the heuristic cannot establish \\
\bottomrule
\end{tabularx}
\caption{A finding record is an editorial question with provenance.}
\label{tab:manual-finding-record}
\end{table}

The author decision is intentionally richer than a checkbox. "Reject" may
mean that the pattern is intentional, that the proposed consequence is wrong,
or that the author will address it in a different paragraph. The reason lets a
later reviewer distinguish a false positive from an unfinished repair. It also
gives the diagnostic designer evidence about which rules waste time.

A model suggestion, when present, is stored as a child of the finding and never
as a replacement for the author's text. The system records whether the author
accepted the wording, adapted it, or wrote an independent repair. This
distinction matters for disclosure and for a later analysis of which assistance
actually reduced work.
\section{Reader form and coding manual}
\label{sec:manual-reader-form}

The reader form should fit on one page plus an optional comment sheet. It opens
with the decision context and then asks:

\begin{enumerate}
\item In one sentence, what decision is the document about?
\item What is the proposed mechanism or object?
\item Which result, equation, or table carries the main support?
\item What condition limits the result?
\item Would you approve the requested next step now? Choose yes, revise, or
      no, and give one reason.
\item Where did you first stop or reread? Give the page or section, a short
      quotation, and the relation you were trying to recover.
\end{enumerate}

The coding manual defines an answer as correct when it identifies the
document-specific target within the declared scope, even if the reader uses
different words. It defines an answer as incomplete when it names a topic but
not the action or mechanism, and as incorrect when it assigns the result to a
different population, comparison, or horizon. These categories are preferable
to asking a model to judge whether the response is eloquent.

For reader $j$, version $v$, and question $q$, let $Z_{jvq}$ take values
$1$ (correct), $0$ (incomplete), and $-1$ (incorrect). The reconstruction index
is reported as a vector of question-level proportions:
\begin{equation}
  \rho_{vq} = \frac{1}{J}\sum_{j=1}^{J}1\{Z_{jvq}=1\},
  \qquad q=1,\ldots,5.
  \label{eq:manual-reader-coding}
\end{equation}
The vector is more informative than its average. A version can improve the
action answer while leaving the mechanism answer unchanged, which tells the
author where a further repair is needed.

The form records confidence and prior exposure separately. Confidence is a
reader observation, not a multiplier on correctness. Prior exposure is a
potential source of contamination, not a reason to discard a skeptical answer.
The evaluator reports both so that a sponsor can see whether the product works
for a new reader or only for someone already familiar with the project.

\section{Evidence and omission register}
\label{sec:manual-evidence-register}

A load-bearing literature row has six linked parts: the source identity, the
technical passage inspected, the source result, its population and task, the
limit, and the local product implication. The record also states what the
source is not allowed to support. For example, an experiment on email
completion can support a bounded claim about task time; it cannot support a
claim about acceptance of an economics proposal.

The omission register uses the following fields:

\begin{table}[H]
\centering
\small
\begin{tabularx}{\textwidth}{@{}p{0.24\textwidth}p{0.27\textwidth}L@{}}
\toprule
Field & Example & Use \\
\midrule
candidate & specialist technical-communication study & Names a plausible
missing source or stream \\
reason considered & Direct reader-comprehension relevance & Explains reviewer risk \\
current status & source blocked, not screened, or adjacent & Prevents silence from
looking like exclusion \\
permitted role & context, direct support, challenge, or none & Limits inference \\
next action & obtain full text, independent screen, or omit with reason & Gives the
review a concrete next step \\
owner and date & named reviewer and deadline & Makes the omission actionable \\
\bottomrule
\end{tabularx}
\caption{A scholarly omission is recorded as a question with a next action.}
\label{tab:manual-omission}
\end{table}

The current evidence run has explicit release blockers for independent
screening, full-text anchors, appraisal, held-out package effectiveness, and
external review. Specialist coverage and DOI identity have passed for the
bounded records recorded in the audit, but those passes do not invalidate the
remaining source-to-claim work or make the survey exhaustive. A later audit can
close a blocker or narrow a sentence; the main route should not silently change
its claim class in response.

\section{Cost worksheet}
\label{sec:manual-cost}

The cost worksheet separates observed inputs from planning assumptions. For
document class $g$, record annual volume $N_g$, incumbent reader hours
$T_{0g}$, proposed reader hours $T_{1g}$, author hours $A_{0g}$ and $A_{1g}$,
reader rate $c_r$, author rate $c_a$, and recurring operating cost $C_o$.
The direct annual value range is
\begin{equation}
  V_g =
  N_g\left[(T_{0g}-T_{1g})c_r+(A_{0g}-A_{1g})c_a\right]-C_o.
  \label{eq:manual-cost-range}
\end{equation}
The worksheet must include a separate line for protected-object incidents.
If the sponsor cannot price an incident, it should report a low, medium, and
high scenario rather than assigning a zero cost.

A worked sensitivity grid varies reader saving and annual volume while holding
the loaded reader cost fixed. The grid is not a forecast. It tells the evaluator
which baseline observations have the largest value of information. If the
break-even point requires an implausibly large volume, the beachhead is wrong
even if the product is technically interesting.

The worksheet also records labor that a software demo hides: intake, source
cleanup, finding investigation, citation review, reader recruitment, and
fixture maintenance. A narrow service with high per-document value may tolerate
more manual work than a broad subscription. The operating plan should price
the actual path observed in the pilot, not an imagined fully automated path.

\section{Decision packet and reproducibility}
\label{sec:manual-packet}

The final packet for a sponsor contains, in order:

\begin{enumerate}
\item the one-sentence decision and the population of documents in scope;
\item the incumbent and proposed workflows, including assistance disclosure;
\item the source and rendered hashes for every compared version;
\item the protected-object comparison and all unresolved items;
\item the reader outcomes, author burden, and missing observations;
\item the evidence status and the sources allowed to support each conclusion;
\item the cost worksheet with sensitivity and non-priced outcomes; and
\item a proceed, revise, or stop recommendation with its next owner.
\end{enumerate}

A packet is reproducible when an independent operator can rebuild the source,
replay the fixtures, locate the protected objects, and recover the same
numerical summaries from the recorded input files. It need not reproduce a
model's stochastic suggestion bit for bit if the model configuration and
output hash are preserved; it must distinguish a stochastic variation from a
source or reader outcome.

The canonical Humanvoice build currently records the top-level source hash,
route hash, rendered hash, page count, extracted word count, and the
evidence-status file. Those fields identify this draft. Release still
requires the project owner and an independent reader, because reproducibility
is evidence about the build and not an endorsement of the prose.

\section{Compact glossary}
\label{sec:manual-glossary}

\begin{description}[style=nextline,leftmargin=36mm,labelindent=0mm]
\item[Protected object] A document element whose change can alter scientific,
financial, or legal meaning and therefore receives explicit comparison.
\item[Reading brief] The declared decision, audience, context, and time
conditions for a reading task.
\item[Finding] A located observation and question that directs human attention.
\item[Abstention] An explicit statement that the current parser, rule, or model
cannot establish correspondence or support.
\item[Incumbent bundle] The actual tools and human work used when Humanvoice is
not present.
\item[Feasibility evidence] Evidence that the path can run and be measured on a
declared case; it is not an efficacy estimate.
\end{description}

The glossary is intentionally short. Domain nouns should be defined where they
enter the argument, and technical field names should remain in this appendix
when they do not help a reader decide whether to fund the next step.
